\subsection{Introduction/prérequis}
\begin{frame}{Le pourquoi du comment}
\begin{block}<only@1-2>{Rappel}
\begin{itemize}
\item<1-> La plongée sous-marine est encadrée par le \emph{Code du Sport}.
C'est une activité pratiquée dans un \og environnement spécifique\fg.
\begin{itemize}
\item[$\Rightarrow$] Elle ne se pratique pas n'importe comment! Il
  y a un cadre légal à respecter.
\end{itemize}
\item<2> La \FFESSM\ est délégataire du ministère \textsl{jeunesse et sports}.
\begin{itemize}
\item[$\Rightarrow$] Elle est l'interlocutrice privilégiée pour modifier et
  faire évoluer la loi.
\end{itemize}
\end{itemize}
\end{block}
%
\begin{block}<only@3>{Niveau 2}
\begin{itemize}
\item PA-20: \textbf{P}longeur \textbf{A}utonome jusqu'à \numprint[m]{20};
\item PE-40: \textbf{P}longeur \textbf{E}ncadré jusqu'à \numprint[m]{40}.
\end{itemize}
\end{block}
\end{frame}

\begin{frame}{Prerequis}
\begin{block}{Toujours}
\begin{itemize}
\item Avoir une licence en cours de validité;
\item être N1 ou équivalent et avoir au minimum quatre plongées en
  milieu naturel;
\item avoir un CACI en cours de validité (moins de un an).
\end{itemize}
\end{block}
\begin{itemize}[<uncover@2->]
\item[PE-40] Être âgé de 14 ans au moins;
\item[PA-20] être âgé de 15 ans au moins;
\item[Niveau~2] être âgé de 15 ans au moins.
\end{itemize}
\begin{alertblock}<uncover@3>{}
Les prérogatives d'évolution en autonomie ne s'exercent qu'à partir de 16 ans.
\end{alertblock}
\end{frame}

\subsection{Officiellement parlant/prérogatives}
\begin{frame}{La FFESSM}{\only<1>{\url{https://ffessm.fr/actualites/la-ffessm-la-connaissez-vous}}\only<2>{\url{https://ffessm.fr/la-ffessm/carte-comites-regionaux}}}
\centering
\only<1>{%
\includegraphics[width=\textwidth]{ffessm_structure}
}%
\only<2>{%
\begin{minipage}{0.44\textwidth}
Organes déconcentrés:
\begin{itemize}
\item comités départementaux,
\item comités régionaux,
\end{itemize}
\end{minipage}\hfill%
\begin{minipage}{0.554\textwidth}
\includegraphics[width=\textwidth]{ffessm_region}
\end{minipage}%
}
\end{frame}

\begin{frame}{Le Code du Sport}{Articles A 322-71 à 101}
\only<1>{Définit les aptitudes PA/PE aux
espaces d'évolution \numprint{12}, \numprint{20}, \numprint{40} et \numprint[m]{60}
(\chref{https://www.legifrance.gouv.fr/codes/article_lc/LEGIARTI000025393959}{Annexe III-14 a (art. A322-77)}).
}
\begin{block}<only@2>{PA-20}
\begin{itemize}
\item Maîtrise des aptitudes PA-12 et PE-20
\item Maîtrise de l'utilisation de l'équipement de ses coéquipiers
\item Maîtrise de sa décompression et du retour en surface à vitesse contrôlée, maintien du palier de sécurité avec parachute de palier
\item Maîtrise d'intervention sur un plongeur en difficulté depuis le fond
\end{itemize}
\end{block}
\begin{block}<only@2>{}
\begin{itemize}
\item[+] Orientation
\end{itemize}
\end{block}
\begin{block}<only@3>{PE-40}
\begin{itemize}
\item Maîtrise des aptitudes PE-20
\item Maîtrise de la vitesse de descente lors de l'immersion
\item Maintien d'un palier avec utilisation d'un parachute
\item Connaissance des signes spécifiques à cette profondeur et maîtrise de la rapidité d'exécution dans les réponses
\item Maîtrise d'une remontée en sécurité en cas de perte de palanquée
\item Intégration à une palanquée guidée à une profondeur de \numprint{20} à \numprint{40}~mètres
\end{itemize}
\end{block}
\begin{alertblock}<only@4>{On oublie trop souvent!}
\begin{itemize}
\item[+] savoir-vivre du plongeur:
  \begin{itemize}
  \item respect de l'environnement,
  \item politesse et aide sur le bateau, avec les autres palanquées, \dots
  \end{itemize}
\end{itemize}
\end{alertblock}
\end{frame}

\begin{frame}{Matériel}
\begin{minipage}{0.4\textwidth}
Nouveautés N2:
\begin{itemize}
\item Ordinateur,
\item parachute,
\item<info@+> deux sorties (octopus)
\end{itemize}
\end{minipage}\hfill
\begin{minipage}{0.5\textwidth}
\includegraphics[width=\textwidth]{N2_prer}
\end{minipage}
\end{frame}

\subsection{Papiers!}
\begin{frame}{Assurance}
\only<1>{%
Trois cas:
\begin{itemize}
\item Non-respect de la loi $\Rightarrow$ pénal, pas d'assurance;
\item dommages avec responsable: \emph{responsabilité civile} (obligatoire);
\item dommages sans responsable: \emph{assurance individuelle} (facultative).
\end{itemize}}%
\includegraphics<2>[width=300pt]{assurances}%
\includegraphics<3>[width=300pt]{assurances_penal}%
\includegraphics<4>[width=300pt]{assurances_civile}%
\includegraphics<5>[width=300pt]{assurances_perso}%
\end{frame}

\begin{frame}{CACI}
\only<1>{%
\begin{minipage}{0.49\textwidth}
\begin{itemize}
\item \myhref{https://plongee.ffessm.fr/uploads/media/docs/0001/13/36282b50927bd779937cb409bb56285f2959c613.pdf}{Le modèle}
se trouve sur le site de la FFESSM.
\item Voir l'\myhref{https://medical.ffessm.fr/certificat-medical-et-pratique-de-l-activite}{article} de la commission médicale
  sur le CACI (en date du \formatdate{24}{02}{2025}).
\end{itemize}%
\end{minipage}
\begin{minipage}{0.49\textwidth}
\includegraphics[width=140pt]{CACI}
\end{minipage}
}
\includegraphics<2>[width=289pt]{CACI_recap}
\end{frame}

\subsection{Fédéralement vôtre}
\begin{frame}{Confédération Mondiale des Activités Subaquatiques}{\url{www.cmas.org/}}
\centering
\begin{itemize}[<only@1>]
\item Fondée en 1959 autour de Jacques-Yves Cousteau;
\item présente dans plus de 100 pays;
\item 181 membres dans le monde (dont 101 fédérations nationales comme FFESSM).
\end{itemize}%
\only<2>{%
\raisebox{-2.5cm}{\includegraphics[height=5cm]{cartes_FFESSM}}\hfill
\begin{tabular}{ll}
\includegraphics[height=1cm,viewport=40 25 210 210,clip]{logoFFESSM} &  \hfill\includegraphics[height=1cm]{logoCMAS}\\\midrule
N1 & plongeur $\bigstar$\\
N2 & plongeur $\bigstar\bigstar$\\
N3 & plongeur $\bigstar\bigstar\bigstar$\\\midrule
E2 & moniteur $\bigstar$\\
E3 & moniteur $\bigstar\bigstar$\\\bottomrule
\end{tabular}\hfill
\raisebox{-2.5cm}{\includegraphics[height=5cm]{cartes_CMAS}}}
\end{frame}

\begin{frame}{Association Nationale des Moniteurs de Plongée}{\url{https://anmp-plongee.com/}}
\centering%
\includegraphics[width=2cm]{logoANMP}
\begin{block}{}
L'ANMP-Guide de la Mer a été créée à Antibes par Daniel Mercier en 1972. Malgré
son nom, l'Anmp n'est pas une association, mais un syndicat professionnel. Elle
agit donc sans but lucratif.

L'ANMP-Guide de la Mer est un syndicat pour faire entendre la voix des
moniteurs professionnels diplômés d'État et promouvoir la profession de
moniteur de plongée.

L'ANMP-Guide de la Mer est une organisation professionnelle pour travailler de
manière efficace et autonome.

L'ANMP-Guide de la Mer crée ses cursus en 1991 et édite ses premières cartes de
certification en 1994 afin de répondre à la demande et aux besoins de ses
adhérents.
\end{block}
\end{frame}

\begin{frame}{Fédération Sportive et Gymnique du Travail}{\url{www.fsgt.org}}
\centering%
\begin{minipage}{1.8cm}
\includegraphics[width=\textwidth]{logoFSGT}
\end{minipage}\hfill
\begin{minipage}{8.6cm}
\begin{block}{}
La Fédération Sportive et Gymnique du Travail (FSGT), doublement agréée sport
et éducation populaire, compte \numprint{200000} pratiquant.e.s,
\numprint{4200} clubs et propose une centaine d'activités physiques, sportives
et artistiques, compétitives et non compétitives. Créée en 1934 au sein du
mouvement sportif ouvrier et dans la lutte contre le fascisme, elle promeut le
droit d'accès au sport de toutes et tous en se donnant comme objectif le
développement de contenus d'activités, de vie associative et de formation
adaptés aux besoins de la population.
\end{block}
\end{minipage}
\end{frame}
